% ============================================================================
\section{Limitations and Conclusion}
\label{sec:conclusion}
% ============================================================================

The study covers one student scale, four tasks, one seed per method, and one
system layout. GPAS is a local rule: its analysis treats micro-batches within
a step as conditionally independent and identically distributed within task,
holds the pre-step preconditioner fixed while
measuring noise, and its attainable variance reduction is bounded by the count
limits. The additive step-time model is fitted to one layout; systems with
parallel critical paths require their own measured time model. The method needs
at least two micro-batches per task per step, so settings with many more tasks
than micro-batches need a task-subsampling variant.
The variance derivation treats fresh micro-batches as conditionally i.i.d.; the
fixed shuffled prompt stream is read without replacement and therefore has a
finite-population covariance correction. Held-out splits and online estimates
measure the realized setting directly.

Each rollout batch is used for one update; staleness caused by reusing rollouts
across updates is outside our scope. The top-$k$ estimator also requires the
teacher to score the student's retained tokens and the sampled token when it
lies outside that set. This is natural when the teacher is available as a
prefill scoring service, but it may
not be available through an API that exposes only the sampled-token score.

GPAS fixes the objective weights and adapts only how many micro-batches each
task contributes, so the allocation changes gradient variance and nothing
else. Its per-step mode follows Neyman allocation in preconditioned coordinates,
and its cost-aware mode also uses measured time; both reuse statistics produced
by MOPD training. \missing{Conclude with the four empirical laws in one
sentence: noise is heterogeneous, drifting, and coordinate dependent; teacher
distance predicts a removable token component; analytic removal and allocation
are complementary; and the overhead ratio bounds the gain per GPU hour.}
